\documentclass[a4paper,11pt]{article}
\usepackage[utf8]{inputenc}
\usepackage[T1]{fontenc}
\usepackage[french]{babel}
\usepackage{lmodern}
\usepackage{amsmath,amssymb,amsthm}
\usepackage{mathrsfs}
\usepackage{enumitem}
\usepackage{multicol}

\setlist[enumerate]{label=\textbf{\textcolor{Couleur8}{\arabic*.}}, left=1em}

% Si vous voulez aussi modifier les listes enumerate imbriquées :
\setlist[enumerate,2]{label=\textcolor{Couleur8}{\alph*)}}
\setlist[enumerate,3]{label=\textcolor{Couleur8}{\roman*)}}
\setlist[enumerate,4]{label=\textcolor{Couleur8}{\Alph*)}}
\usepackage{graphicx}
\usepackage{xcolor}
\usepackage{tcolorbox}
\usepackage{geometry}
\usepackage{fancyhdr}
\usepackage{titlesec}
\usepackage{cancel}
\usepackage{ifthen}
\usepackage{tikz}
\usetikzlibrary{arrows.meta}
\usepackage{tkz-tab}
\usepackage{eurosym}

% OPTION PROF/ÉLÈVE
% Décommentez UNE des deux lignes suivantes :
\newboolean{prof}\setboolean{prof}{true}   % Version professeur (avec corrections des devoirs)
%\newboolean{prof}\setboolean{prof}{false}  % Version élève (sans corrections des devoirs)

\definecolor{Couleur1}{HTML}{4A5D7A} %Th
\definecolor{Couleur2}{HTML}{A5B5D6} %Prop
\definecolor{Couleur3}{HTML}{A8C68A} %Méthode
\definecolor{Couleur6}{HTML}{8A9CC1} %Def
\definecolor{Couleur7}{HTML}{E6B459} %Rq Voc
\definecolor{Couleur8}{HTML}{D36740} %Titres
\definecolor{correction}{HTML}{D36740} % Couleur pour les corrections
\definecolor{coopmathscorr}{HTML}{F15929} % Couleur corrections coopmaths

% Configuration des marges
\geometry{
	top=2cm,
	bottom=2cm,
	inner=1.5cm,
	outer=1.5cm,
}

% Police
\renewcommand{\familydefault}{\sfdefault}

% Compteurs
\newcounter{seancenum}
\newcounter{seancenumD}
\setcounter{seancenum}{1}
\setcounter{seancenumD}{1}

% Configuration des en-têtes et pieds de page
\pagestyle{fancy}
\fancyhf{}
\ifthenelse{\boolean{prof}}{
    \fancyhead[L]{\textcolor{Couleur8}{\textbf{Corrections Automatismes - Classe de seconde - Version Professeur}}}
}{
    \fancyhead[L]{\textcolor{Couleur8}{\textbf{Corrections Automatismes - Séance 4 - Classe de seconde}}}
}
\fancyhead[R]{\textcolor{Couleur8}{\textbf{2025-2026}}}
\fancyfoot[L]{\textcolor{Couleur8}{Mme Bertail et Mme Pinto}}
\fancyfoot[R]{\textcolor{Couleur8}{\thepage}}
\renewcommand{\headrulewidth}{1pt}
\renewcommand{\headrule}{\hbox to\headwidth{\color{Couleur8}\leaders\hrule height \headrulewidth\hfill}}

% Environnements personnalisés
\newtcolorbox{seance}[1]{
    colback=Couleur6!10,
    colframe=Couleur1!65,
    fonttitle=\bfseries\large,
    title=Correction Séance \theseancenum #1,
    left=5mm,
    right=5mm,
    top=3mm,
    bottom=3mm,
    arc=0mm,
    after=\stepcounter{seancenum}
}

\newtcolorbox{seanceD}[1]{
    colback=Couleur6!5,
    colframe=Couleur6!45,
    fonttitle=\bfseries\large,
    title=Correction Devoirs - Séance \theseancenumD #1,
    left=5mm,
    right=5mm,
    top=3mm,
    bottom=3mm,
    arc=0mm,
    after=\stepcounter{seancenumD}
}

\begin{document}

\setcounter{seancenum}{4}
\newpage
\begin{seance}{} %4

\begin{enumerate}
\item La seule égalité vraie est : ${\color{correction}\boldsymbol{\dfrac{15^{-7}}{15^{11}}=15^{-18}}}$.
En effet, $\dfrac{15^{-7}}{15^{11}}=15^{-7-11}=15^{-18}$

Concernant les autres propositions :
$20\times \dfrac{1}{20^{5}}=20^{-4}\neq 20^{4}$
$\left(40^{-5}\right)^4=40^{-20}\neq 40^{-1}$
$7^{-5}\times 5^{-5}=35^{-5}\neq 35^{-10}$

\item $N=\dfrac{6^5}{2^3}=\dfrac{2^5\times 3^5 }{2^3}=3^5\times 2^{2}=6^2\times 3^{3}={\color{correction}\boldsymbol{27\times 6^{2}}}$

\item  On utilise la formule $\dfrac{a^n}{b^n}=\left(\dfrac{a}{b}\right)^{n}$ avec
        $a=6$,  $b=3$ et $n=-8$.\\
        $\dfrac{6^{-8}}{3^{-8}}=\left(\dfrac{6}{3}\right)^{-8}={\color{correction}\boldsymbol{2^{-8}}}$
        


\item On applique la propriété des puissances de puissances d'un réel : \\
    Soit $n\in \mathbb{N}$, et $p \in \mathbb{N}$, on a : \\
     $\left(a^{n}\right)^{p}=a^{np}$\\
    $\begin{aligned}\left(3^{n}\right)^{3}&=\left(3\right)^{3n}\\
    &=\left(3^{3}\right)^{n}\\
    &=27^{n}
    \end{aligned}$\\La bonne réponse est la réponse {\bfseries \color{correction}C}.


\item De la relation $6a+5y=4$, on déduit en ajoutant $-5y$ dans chaque membre :
$6a=4-5y$
Puis, en divisant par $6$, on obtient : ${\color{correction}\boldsymbol{a=\dfrac{4-5y}{6}}}$.






\item La réunion $]-3\,;\,3[\,\cup \,] -6\,;\,5]$ = ${\color{correction}\boldsymbol{]-6\,;\,5]}}$ car l'intervalle $] -6\,;\,5]$ contient entièrement l'intervalle $]-3\,;\,3[$.

\end{enumerate}

\end{seance}

\end{document}